\clearpage
\SourceLeaf{105}%

\begin{center}
\large A MOST ELEGANT THEOREM.
\end{center}

\bigskip

\noindent\hspace*{\parindent}Let

\[ 1+\tfrac{1}{2}\cdot\tfrac{1}{2}xx+\tfrac{1}{2}\cdot\tfrac{1}{2}\cdot\tfrac{3}{4}\cdot\tfrac{3}{4}x^4+\tfrac{1}{2}\cdot\tfrac{1}{2}\cdot\tfrac{3}{4}\cdot\tfrac{3}{4}\cdot\tfrac{5}{6}\cdot\tfrac{5}{6}x^6+\text{etc.} = fx, \]

\[ \tfrac{1}{2}x+\tfrac{1}{2}\cdot\tfrac{3}{2}\cdot\tfrac{3}{4}x^3+\tfrac{1}{2}\cdot\tfrac{3}{2}\cdot\tfrac{3}{4}\cdot\tfrac{5}{4}\cdot\tfrac{5}{6}x^5+\text{etc.} = f'x, \]

\noindent then we shall have

\[ \sin\varphi\, f\sin\varphi\, f'\cos\varphi + \cos\varphi\, f\cos\varphi\, f'\sin\varphi = \frac{2}{\pi\sin\varphi\cos\varphi}. \]

\bigskip

\bigskip

\begin{center}
REMARKS ON »THEOREMA ELEGANTISSIMUM«.
\end{center}

\noindent\hspace*{\parindent}\textsc{Gauss} wrote this theorem on the last page of his own copy of »\textsc{Euler}, Methodus inveniendi curvas Maximi Minimive proprietate gaudentes«. The theorem was included here because of the elegant form of the concluding formula, which evidently struck \textsc{Gauss} himself. In substance, this is one of the simplest differential relations in the theory of the arithmetic-geometric mean, as can be shown with the aid of a few formulas from Art.\,13 and 14 of »Arithmetisch-geometrischen Mittels« (cf. ges. Werke III p.\,379 ff.). For $fx$ proves to be the \textsc{Gauss}ian series $F(\tfrac{1}{2},\tfrac{1}{2},1;x^2)$, so that, using the last two formulas on p.\,381 l.\,c., we may put:

\[ x=\sin\varphi=\frac{b}{a}, \qquad \sqrt{1-x^2}=\cos\varphi=\frac{c}{a}, \]

\[ f\sin\varphi=\frac{a}{M( a,c)}, \qquad f\cos\varphi=\frac{a}{M( a,b)}. \]

\noindent After a simple intermediate calculation, the theorem stated here then takes the form:

\[ \frac{b^2\mathord{.} d\log\dfrac{a}{M( a,b)}}{d\log\dfrac{c}{a}} - \frac{c^2\mathord{.} d\log\dfrac{a}{M( a,c)}}{d\log\dfrac{a}{b}} = \frac{2}{\pi}\,M( a,b)\,M( a,c), \]

\noindent a formula that is a simple consequence of the first equations on p.\,380 l.\,c.

\hfill \textsc{Fricke}.

\bigskip

